\section{Introduction}
\label{sec:intro}
Radar underpins applications ranging from air defence and space situational awareness to weather forecasting and automotive safety \cite{palmerHorusFullyDigital2023, engelsAutomotiveRadarSignal2021, richardsPreviewBasicRadar2022}. At its core lies the signal processing chain that turns raw echoes into information an operator can use. This chain typically involves pulse compression, doppler processing, beamforming as well as detection algorithm such as Constant False Alarm Rate (CFAR). For Synthetic Aperture Radar (SAR), image formation is also included in the signal processing chain. Over the past two decades, the computational demand for this chain has grown considerably. This is largely down to two reasons. First, wider bandwidths and element-level digitisation deliver far larger volumes of data to the processor \cite{fultonDigitalPhasedArrays2016, uekiDataChallengeFully2025}. Secondly, algorithms have shifted from fixed to adaptive forms whose cost grows steeply with problem size \cite{bensonGPUbasedSpacetimeAdaptive2013a, fultonDigitalPhasedArrays2016}. At the same time, more of the receive chain is moving from bespoke hardware and firmware into software. This software-defined radar, analogous to software-defined radio, enables configurable waveforms, modes, and processing change without redesigning the hardware, but also adds a computational load on the processors.

Meeting these demands within real-time deadlines has driven radar processing towards heterogeneous computing architectures that pair a host processor with dedicated accelerators \cite{mittalSurveyCPUGPUHeterogeneous2015}. Field Programmable Gate Arrays (FPGAs) have filled this role and remain favoured for communication heavy front-end stages \cite{yuImplementationRealTimePhased2016}. However, they are costly to develop and difficult to reconfigure, which sits poorly with a software defined chain \cite{perdanaSurveyGraphicsProcessing2017}. Graphics Processing Units (GPUs) have emerged as an attractive alternative by providing commercial off-the-shelf (COTS) devices with high arithmetic throughput and high memory bandwidth alongside a mature software stack. Initial reported speedups for radar workloads range from tens to hundreds of times \cite{perdanaSurveyGraphicsProcessing2017, zhaoGPUAcceleratedSignalProcessing2023, zhangAcceleratingSpaceborneSAR2016a}.

Despite these results, the published work explains little about why GPU implementations perform as they do.
Stages such as the corner turn, envelope detection and windowing perform little arithmetic per byte moved,
so memory and data movement bound them rather than computation
\cite{williamsRooflineInsightfulVisual2009a, venterSOFTWAREDEFINEDPULSEDOPPLERRADAR, denhamWeatherRadarData2018a},
yet studies rarely identify which resource bounds each stage.
Headline speed-ups are often measured against serial CPU baselines
\cite{zhaoGPUAcceleratedSignalProcessing2023, zhangAcceleratingSpaceborneSAR2016a, denhamWeatherRadarData2018a, perdanaRadarSignalProcessing2019, yangGPUOrientedDesignsConstant2022a},
which overstates the GPU's advantage \cite{leeDebunking100XGPU2010b};
transfers are timed inconsistently \cite{venterSOFTWAREDEFINEDPULSEDOPPLERRADAR, greggWhereDataWhy2011a};
timings rarely report spread \cite{hoeflerScientificBenchmarkingParallel2015a};
and execution features such as CUDA Graphs, designed for repeated fixed-shape workloads like a radar pipeline,
go unused \cite{ekelundBoostingPerformanceIterative2025a}.
It therefore remains unclear which implementation strategies, rather than which hardware, drive the reported gains.

This review addresses the gap from a computational rather than algorithmic perspective. The main review question that is asked is \emph{which implementation strategies, rather than which hardware, explain the performance reported for GPU radar signal processing, and what do current studies leave unmeasured?} Studies were identified by citation tracking from seed papers, with a database search used to check coverage, and each study was appraised against six reporting criteria drawn from established benchmarking practice \cite{hoeflerScientificBenchmarkingParallel2015a, greggWhereDataWhy2011a}. The review covers the processing from pulse compression to CFAR detection, where data volumes are largest, with SAR image formation as a compute-intensive contrast. The paper aims to make three contributions (i)~ a per-stage characterisation of the radar chain as a computational workload, identifying which stages are memory-bound and which compute-bound (Section ~\ref{sec:workload}); (ii)~ an appraisal of how GPU performance is measured and reported, including how data movement is accounted for (Sections~\ref{sec:reporting} and ~\ref{sec:data}); and (iii)~ an account of the implementation strategies likely to matter most but untested on current hardware, from which research questions are derived (Sections~\ref{sec:execution} and ~\ref{sec:gaps}). Section ~\ref{sec:background} introduces the radar receive chain and the relevant GPU features, Section ~\ref{sec:method} outlines the review method, and Section ~\ref{sec:conclusion} concludes by drawing together the grand challenge.
